\subsection{Приказ ФСТЭК России № 21}

Приказ ФСТЭК России от 18.02.2013 № 21 «Об утверждении Состава и содержания организационных и технических мер по обеспечению безопасности персональных данных при их обработке в информационных системах персональных данных» \cite{FSTEK21} является одним из ключевых подзаконных актов, конкретизирующих требования к защите персональных данных в информационных системах персональных данных. Данный нормативный документ устанавливает состав мер, необходимых для обеспечения безопасности персональных данных, и определяет подходы к их реализации в зависимости от уровня защищённости информационной системы.

В соответствии с приказом № 21 все меры защиты подразделяются на организационные и технические и охватывают полный жизненный цикл обработки персональных данных. К основным группам мер относятся:

\begin{itemize}
    \item \textbf{Идентификация и аутентификация субъектов доступа} — обеспечение однозначного установления личности пользователей и процессов, имеющих доступ к информационной системе, включая применение механизмов аутентификации и управление учётными записями.
    \item \textbf{Управление доступом} — разграничение прав доступа пользователей и процессов к информационным ресурсам в соответствии с их полномочиями, реализация принципа минимально необходимых привилегий, а также контроль доступа к данным и функциям системы.
    \item \textbf{Ограничение программной среды} — предотвращение запуска несанкционированного программного обеспечения, контроль состава и целостности используемых программных компонентов.
    \item \textbf{Защита машинных носителей информации} — обеспечение контроля доступа к носителям информации, предотвращение несанкционированного копирования и утечки данных.
    \item \textbf{Регистрация событий безопасности} — ведение журналов событий, связанных с доступом к системе и действиями пользователей, включая фиксацию попыток несанкционированного доступа, изменений конфигурации и иных значимых событий.
    \item \textbf{Антивирусная защита} — применение средств защиты от вредоносного программного обеспечения и контроль их актуальности.
    \item \textbf{Обнаружение вторжений} — выявление попыток несанкционированного доступа и иных инцидентов безопасности на основе анализа событий и поведения пользователей.
    \item \textbf{Контроль целостности} — обеспечение неизменности программного обеспечения и данных, включая контроль целостности журналов регистрации событий.
    \item \textbf{Обеспечение доступности} — защита системы от отказов и сбоев, а также обеспечение возможности восстановления работоспособности.
    \item \textbf{Защита среды виртуализации} — обеспечение безопасности виртуальных инфраструктур, включая изоляцию компонентов и контроль взаимодействия между ними.
\end{itemize}

Перечисленные меры формируют комплексную систему защиты информации и направлены на предотвращение, выявление и реагирование на инциденты информационной безопасности.

Для темы ВКР данный приказ имеет особое значение, поскольку разрабатываемое средство основано на анализе audit-логов Kubernetes-кластера, то есть на сборе, регистрации и последующей обработке событий, связанных с действиями пользователей и сервисных учётных записей. Это напрямую соответствует требованиям по регистрации событий безопасности и обнаружению вторжений, установленным приказом № 21.

В контексте Kubernetes-кластера указанные меры могут быть реализованы следующим образом:
\begin{itemize}
    \item механизмы идентификации и аутентификации — через средства аутентификации Kubernetes;
    \item управление доступом — с использованием ролевой модели доступа (RBAC);
    \item регистрация событий — посредством audit-логирования kube-apiserver;
    \item обнаружение вторжений — с использованием разработанной системы анализа аномалий;
    \item контроль целостности — через защиту журналов и централизованные системы хранения логов.
\end{itemize}